\documentclass[12pt,a4paper]{article}
% Preambule commun aux comptes rendus CR_*.tex.

\usepackage{array}
\newcolumntype{C}[1]{>{\centering\arraybackslash}m{#1}}
\usepackage{multirow}
\usepackage{geometry}
\geometry{a4paper, margin=1in}
\usepackage{graphicx}
\usepackage{float}
\usepackage{tikz}
\usetikzlibrary{arrows.meta,positioning}
\usepackage{pgfgantt}
\usepackage{pdflscape}
\usepackage{enumitem}
\usepackage{booktabs}
\usepackage{longtable}
\usepackage{ragged2e}
\usepackage{tabularx}
\usepackage{xparse}
\usepackage{ltablex}
\usepackage{xltabular}
\usepackage{subcaption}

\newcolumntype{Y}{>{\RaggedRight\arraybackslash}X}

\newcommand{\StyledTableSetup}{%
  \footnotesize
  \renewcommand{\arraystretch}{1.2}%
  \setlength{\tabcolsep}{4pt}%
}


% ============================================================
% IHEDN COVER TEMPLATE (XeLaTeX)
% ============================================================
% Compilation (from project root):
%   latexmk -pdf main.tex
%   LATEXMK_SHELL_ESCAPE=0 latexmk -pdf main.tex
%
% Optional environment controls from latexmkrc:
%   LATEXMK_ENGINE=xelatex        (default/enforced)
%   LATEXMK_SHELL_ESCAPE=1|0      (default 1, needed for SVG conversion)
%
% Required package for this template:
%   \usepackage{ihedn-cover}
%
% Note: ihedn-cover already loads needed internal packages
% (fontspec, graphicx, tikz, ragged2e, optional svg).
% ============================================================

% ----------------------------------------------------------------
% Package options (documented)
% ----------------------------------------------------------------
% Geometry scope:
%   apply-geometry            -> force A4 + zero margins while cover is built
%   no-apply-geometry         -> do not alter host geometry (default)
%
% Typesetting freeze:
%   global-typesetting        -> apply freeze globally (intrusive)
%   no-global-typesetting     -> no global freeze (default)
%
% Small-caps handling:
%   local-smallcaps-fix       -> \textsc -> \MakeUppercase only inside cover (default)
%   no-local-smallcaps-fix    -> disable local fix
%   global-smallcaps-fix      -> global override (intrusive)
%   no-global-smallcaps-fix   -> disable global fix (default)
%
% Main font handling:
%   global-mainfont           -> set Times New Roman globally (default)
%   local-mainfont            -> set Times only during cover rendering
%   no-mainfont-override      -> keep host document main font
%
% SVG support:
%   svg-support               -> enable SVG logos (default, needs shell-escape)
%   no-svg-support            -> disable SVG path handling (pdf/png only)
%
% Debug overlay:
%   debug                     -> show mm grid + bounding boxes on cover
%   no-debug                  -> hide debug overlay (default)
% ----------------------------------------------------------------
\usepackage[
  no-apply-geometry,
  no-global-typesetting,
  global-smallcaps-fix,
  local-smallcaps-fix,
  global-mainfont,
  svg-support,
  no-debug
]{ihedn-cover}

% ----------------------------------------------------------------
% Bibliography stack (BibLaTeX/Biber, style from subtree vendor/)
% ----------------------------------------------------------------
\usepackage{polyglossia}
\setdefaultlanguage{french}
\makeatletter
\g@addto@macro\captionsfrench{%
  \renewcommand{\tablename}{Tableau}%
  \renewcommand{\figurename}{Figure}%
}
\makeatother
\usepackage{csquotes}
\usepackage{xurl}
\usepackage{hyperref}
\makeatletter
\let\ihedn@orig@input@path\input@path
\def\input@path{{vendor/biblatex-sciences-po-fr/style/}}
\makeatother
\usepackage[
  backend=biber,
  style=sciencespo,
  hyperref=true,
  autolang=other
]{biblatex}
\AtBeginBibliography{\sloppy\setlength{\emergencystretch}{3em}}
\makeatletter
\let\input@path\ihedn@orig@input@path
\makeatother
\addbibresource{bibliographies/bib/rapport.bib}

% ----------------------------------------------------------------
% tcolorbox
% ----------------------------------------------------------------
\usepackage[most]{tcolorbox}
\tcbuselibrary{breakable}

% Example: use Arial for the full report + cover.
% Uncomment and switch package option to no-mainfont-override.
% \setmainfont{Arial}[
%   UprightFont    = *,
%   BoldFont       = * Bold,
%   ItalicFont     = * Italic,
%   BoldItalicFont = * Bold Italic
% ]

%==================================================
% Liste des propositions
%==================================================

\makeatletter

\newcommand{\listpropositionname}{Liste des propositions}

\newcommand{\listofpropositions}{
  \section*{\listpropositionname}
  \addcontentsline{toc}{section}{\listpropositionname}
  \@starttoc{lop}
}

\newcommand*\l@proposition{\@dottedtocline{1}{1.5em}{2.3em}}

\makeatother

%==================================================
% Compteur des propositions
%==================================================

\newcounter{proposition}

%==================================================
% Style graphique des propositions
%==================================================

\newtcolorbox{propositionbox}[2][]{
  lower separated=false,
  colback=white!80!gray,
  colframe=white!20!black,
  fonttitle=\bfseries,
  colbacktitle=white!30!gray,
  coltitle=black,
  enhanced,
  sharp corners,
  attach boxed title to top left={
    xshift=0.5cm,
    yshift=-2mm
  },
  title=#2,
  #1
}

%==================================================
% Environnement proposition
%==================================================

\newenvironment{proposition}[1]
{
  \refstepcounter{proposition}

  \addcontentsline{lop}{proposition}{%
    \protect\numberline{\theproposition}#1%
  }

  \begin{propositionbox}
  {Proposition~\theproposition~: #1}
}
{
  \end{propositionbox}
}

% Renommage Liste des figures en Table des figures
\addto\captionsfrench{%
  \renewcommand{\listfigurename}{Table des figures}
}

\newcommand{\heure}[2]{{\NoAutoSpacing #1:#2}}

\DeclareFieldFormat[misc]{title}{#1}

% Activer ce package avec l'option none retire la césure.
%\usepackage[none]{hyphenat}

\usepackage{adjustbox}
\usepackage{pdfpages}

\makeatletter
\newcommand{\IhednSetPdfPageCount}[2]{%
  \begingroup
  \edef\ihedn@pdfpath{#2}%
  \ifXeTeX
    \global#1=\XeTeXpdfpagecount "\ihedn@pdfpath" %
  \else
    \ifLuaTeX
      \saveimageresource{\ihedn@pdfpath}%
      \global#1=\lastsavedimageresourcepages\relax
    \else
      \pdfximage{\ihedn@pdfpath}%
      \global#1=\pdflastximagepages\relax
    \fi
  \fi
  \endgroup
}
\makeatother

\newcount\IhednIncludedPdfPageCount
\newcommand{\IhednRenderPdfPagesInBoxes}[1]{%
  \IhednSetPdfPageCount{\IhednIncludedPdfPageCount}{#1}%
  \foreach \IhednIncludedPdfPage in {1,...,\the\IhednIncludedPdfPageCount}{%
    \begin{tcolorbox}[
      colback=gray!5,
      colframe=black,
      boxrule=0.5pt,
      arc=3mm
    ]
    \centering
    \includegraphics[page=\IhednIncludedPdfPage,width=0.95\textwidth]{#1}
    \end{tcolorbox}
    \ifnum\IhednIncludedPdfPage<\IhednIncludedPdfPageCount
      \clearpage
    \fi
  }%
}

\usepackage{tablefootnote}


\begin{document}

% ----------------------------------------------------------------
% Cover keys (all available keys are shown below)
% ----------------------------------------------------------------
% Header block (top-right), logo, title, report lines, subtitle, committee.
\PageDeGardeIHEDN{
    header_1={Union-IHEDN},
    header_2={Association des auditeurs IHEDN},
    header_3={région Paris / Île de France},
    date={le 26 janvier 2026},
    logo_path={Figures/logos_IHEDN/Logo_AR_16_PARIS_IDF.jpg},
    title={Crises majeures : quel rôle pour les citoyens engagés et les associations IHEDN ?},
    title_max_lines=2,
    title_fontsize={26.000pt},
    title_baselineskip={28.667pt},
    reportline_1={Rapport du Comité d'étude de l'Association régionale des auditeurs},
    reportline_2={de l'IHEDN région Paris / Île de France (AR 16)},
    subtitle={},
    committee={Comité d'étude, d'octobre 2025 à juin 2026, composé de : Mme Valérie Arlé-Faivre, \textit{présidente} ; M. Aurélien Duvignac-Rosa, \textit{rapporteur} ; M. Pascal Bayot Duponchel ; M. Jean-Jacques Cleynen ; Mme Valérie Cowan ; M. Alain Fontan ; M. Pierre-Marie Giard ; Mme Muriel Joyeux ; M. Bruno Leray ; M. Yves-Henri Renhas ; M. Marc Roure ; Mme Isabelle de Segonzac ; M. Gérard Turck ; Mme Marie Danielle Vasquez-Duchêne.}
}

\tableofcontents

\newpage

\section{Introduction}

Les citoyens engagés et les associations régionales de l'IHEDN (Institut des hautes études de défense nationale) peuvent jouer un rôle opérationnel essentiel à toutes les phases du cycle de crise — prévention, gestion et résilience — en complément de l'action de l'État et des collectivités territoriales.

\section{En amont des crises : un rôle clé de prévention et d'anticipation}

\subsection{Diffusion de la culture de sécurité et de résilience}

Les associations régionales de l'IHEDN sont particulièrement bien placées pour :

\begin{itemize}
    \item Sensibiliser les citoyens, élus et acteurs économiques aux risques majeurs (naturels, sanitaires, cyber, énergétiques, informationnels, géopolitiques) ;
    \item Promouvoir l'esprit de défense, la culture de sécurité nationale et la compréhension des enjeux stratégiques ;
    \item Contribuer à l'appropriation des plans de prévention (PCS, PPR, ORSEC) par le grand public.
\end{itemize}

Citoyens engagés : relais locaux, formateurs bénévoles, animateurs de conférences ou d'ateliers de sensibilisation.

\subsection{Veille territoriale et remontée d'informations}

Grâce à leur ancrage local et à la diversité de leurs membres :
\begin{itemize}
    \item Les associations IHEDN peuvent participer à une veille informelle (signaux faibles, vulnérabilités territoriales, tensions sociales, dépendances critiques) ;
    \item Elles peuvent servir de pont de dialogue entre société civile, entreprises, forces de sécurité et autorités locales.
\end{itemize}

\section{Pendant la crise : un appui à la gestion et à la cohésion}

\subsection{Soutien à la continuité sociale et institutionnelle}

Même sans rôle décisionnel, les réseaux IHEDN peuvent :
\begin{itemize}
    \item Contribuer à la continuité des échanges entre acteurs publics et privés ;
    \item Faciliter la compréhension des décisions publiques et limiter les incompréhensions ou défiances ;
    \item Renforcer la cohésion nationale et territoriale dans les moments de tension.
\end{itemize}

Les citoyens engagés formés aux enjeux de crise jouent ici un rôle de médiateurs informés.

\subsection{Appui aux autorités et aux dispositifs existants}

En complément des dispositifs officiels :
\begin{itemize}
    \item Mobilisation de compétences spécifiques (logistique, communication de crise, cybersécurité, santé, gestion des organisations) ;
    \item Participation encadrée à des cellules de réflexion, de soutien ou de retour d'expérience à chaud ;
    \item Relais d'information fiable pour lutter contre la désinformation et les rumeurs.
\end{itemize}

\section{Après la crise : renforcer la résilience et le retour à l'équilibre}

\subsection{Capitalisation et retour d'expérience (RETEX)}

Les associations IHEDN peuvent :
\begin{itemize}
    \item Analyser les enseignements stratégiques, organisationnels et humains de la crise ;
    \item Produire des contributions indépendantes utiles aux décideurs ;
    \item Favoriser la diffusion des bonnes pratiques à l'échelle régionale et nationale.
\end{itemize}

\subsection{Reconstruction du lien social et de la confiance}

Les citoyens engagés et les réseaux associatifs :
\begin{itemize}
    \item Participent à la reconstruction du tissu social ;
    \item Accompagnent les acteurs locaux dans l'adaptation post-crise ;
    \item Contribuent à restaurer la confiance entre citoyens et institutions.
\end{itemize}

\section{Une valeur ajoutée spécifique des associations IHEDN}

Les associations régionales de l'IHEDN apportent :
\begin{itemize}
    \item Une approche transversale (civile, militaire, économique, académique) ;
    \item Une culture du temps long et de la complexité stratégique ;
    \item Un positionnement non partisan, facilitant le dialogue ;
    \item Un réseau structuré, crédible et reconnu par les pouvoirs publics.
\end{itemize}

\section{Conclusion}

Les citoyens engagés et les associations régionales de l'IHEDN ne se substituent pas aux autorités, mais constituent une force d'appoint stratégique et opérationnelle.

Par la préparation des esprits, la mobilisation des compétences, la cohésion sociale et la capitalisation des enseignements, ils contribuent directement à la robustesse, à l'agilité et à la résilience de la Nation face aux crises majeures.

\end{document}